\section{Introduction}\label{sec:intro}

In September 2026 the Italian Parliament gave final approval to a law delegating to the government the reintroduction of nuclear power \citep{Delega2026}. The law does not authorise any plant and does not select sites; the national energy and climate plan foresees 8--16\,GW of nuclear capacity by 2050, equal to 11--22\% of electricity demand. How much nuclear capacity a zero-emission Italian power system should contain, and under which conditions, is therefore a practical question.

The public debate in Italy relies mostly on the levelised cost of electricity (LCOE). In the IEA World Energy Outlook 2024, utility-scale solar in the European Union in 2030 costs about 35\,USD/MWh against about 140\,USD/MWh for nuclear \citep{IEA2024WEO}; the same comparison is common in the Italian debate \citep{Pagella2026}. Industry reports such as \citet{Lazard2025} and \citet{Fraunhofer2024} reach similar conclusions. The LCOE, however, divides the cost of a plant by the energy it produces regardless of when it produces it. In Italy solar output in January is one third of July output (Section~\ref{sec:monthly}), so a solar MWh produced at noon in July and a nuclear MWh produced on a January evening are not the same product. System studies address this. \citet{Sepulveda2018}, analysing nearly 1{,}000 cases for US systems, find that firm low-carbon resources reduce the cost of electricity by 10--62\% in fully decarbonised cases, which is an increase of 11--163\% in the cost without them (our conversion); the \citet{NEA2019} estimates integration costs of variable renewables rising from about 8 to about 50\,USD/MWh between 10\% and 75\% penetration.

For Italy, \citet{Lombardi2020} find a near-optimal 2050 system with 144\,GW of solar, 76\,GW of wind and synthetic methane for back-up, without nuclear. \citet{Mastrantuono2026} find that new nuclear lowers total system cost by up to 8 billion euros, roughly halves imports and reduces storage needs by 42\%, curtailment by 35\% and land use by 25\%; their model, however, caps nuclear at 8\,GWe, and the cap binds in the scenarios with stable costs, so the optimal nuclear share is not identified. An analysis by Politecnico di Milano and Universit\`a dell'Insubria, reported in the trade press, finds up to 15\,GW of nuclear at 5{,}500\,EUR/kW and none at 8{,}300\,EUR/kW with ten-year construction; the model has not been published \citep{PoliMi2026}.

This paper contributes in five ways. First, it estimates the cost-optimal amount of nuclear in a zero-emission Italian system without exogenous caps, for a 2040 demand that includes heat pumps and electric vehicles, with a model at three-hour resolution in which one fleet must serve four weather years, as a planner would require, and with an explicit adequacy reserve. Second, it prices renewables and storage from the prices actually awarded in Italian auctions, rather than from engineering estimates, and converts these nominal, partly indexed contract prices into real costs over the life of the plants, so that they are comparable with a nuclear cost computed in real terms over its own life. Solar and wind follow finite supply curves, because ground-mounted modules are banned on agricultural land and wind sites are limited, so that land enters the cost of the system; existing plants enter at operating cost only. Third, it makes the dependence of the result on the three variables that drive nuclear cost, construction cost, cost of capital and construction time, explicit by summarising them in a single quantity, the annualised fixed cost per kW (fixN), so that the result is one curve and each cost case is a point on it. Fourth, it includes every technology that can produce, store or shift electricity in Italy under one cost rule, with first-plant and target cases for technologies without commercial plants. Fifth, it treats imports as a dependence choice with a price: the central rule is net self-sufficiency, with exports allowed and no foreign help in the critical weeks, and the cost saved by allowing net imports is reported separately, so that a plan is not made to look cheaper by counting on plants that other countries build. The model has one node and no ramp constraints: the aim is to identify the variables that matter and to make the result reproducible, not to represent every feature of the Italian system. It describes a regime system, with 2040 demand and all fleets already built, not the situation in 2040. The result also supersedes an earlier version of this analysis that gave 38.5\,GW and 59\% of demand, in a model without imports, hydrogen, bioenergy, geothermal and the other technologies now included; Section~\ref{sec:bridge} explains the difference step by step. Section~\ref{sec:methods} describes the model and the data, Section~\ref{sec:results} presents the results, Section~\ref{sec:discussion} discusses them together with supply security, industrial capability and political feasibility, and Section~\ref{sec:conclusions} concludes.
